\section{Problem description}
\label{sec:problem}

A plant assembles products supplied by a single supplier over a finite
horizon. Every product travels in exactly one type of RTI: a rigid
container, \emph{erected} while it carries product and \emph{folded}
when empty, so one physical container occupies two different volumes
depending on which way it is moving. The supplier packs its production
into erected RTIs and ships them to the plant; the plant consumes the
product, folds the containers and returns them. Both legs are run with
identical trailers, every dispatch costs the same whatever it carries,
and a trailer may be dispatched only if it is full. A dispatch arrives
$lt$ periods after it leaves and is available at neither end in
between, and each node must end every period holding enough stock to
cover the next periods of its own activity. What has to be decided is,
per direction and period, how many trucks go, what each carries and
what consequently sits in stock; the cost is the number of dispatches
plus the inventory held.

Two features separate this from a transport-and-inventory model. The
two flows are coupled through the containers themselves --- the plant
can return only empties it has emptied and the supplier can fill only
containers it has received --- so a shortage on the return leg
propagates into the delivery leg one lead time later. And the
admissibility of a truck load is a three-dimensional packing condition
rather than a capacity inequality: a load whose volume and weight both
fit may still have no arrangement inside the trailer. Loading is
represented through \emph{stacks} of whole homogeneous layers rather
than through the placement of individual RTIs, which makes the stacks
self-supporting, so no vertical-stability constraint is needed.
